\section{Bối cảnh hình thành đề tài}
AI tạo sinh đã thay đổi mạnh mẽ cách con người sản xuất nội dung số. Trước đây, việc tạo một bức ảnh chân thực hoặc một đoạn nội dung thuyết phục thường cần thiết bị chuyên dụng, kỹ năng thiết kế và thời gian biên tập. Hiện nay, người dùng phổ thông có thể tạo ảnh, bài viết, video ngắn, giọng nói hoặc mô phỏng khuôn mặt chỉ bằng một lời nhắc ngắn. Điều này đem lại nhiều giá trị tích cực cho giáo dục, sáng tạo nội dung, truyền thông và nghiên cứu.

Tuy nhiên, mặt trái của xu hướng này là sự gia tăng của nội dung không rõ nguồn gốc. Một bức ảnh AI có thể được dùng để minh họa sự kiện không có thật; một đoạn văn AI có thể được dùng để tạo bình luận giả; một video deepfake có thể gây ảnh hưởng đến uy tín cá nhân hoặc tổ chức. Khi nội dung số trở nên khó phân biệt bằng mắt thường, xã hội cần các công cụ hỗ trợ kiểm chứng để giảm rủi ro tiếp nhận sai thông tin.

SourceVerify ra đời từ nhu cầu đó. Hệ thống không cố gắng thay thế chuyên gia giám định số, mà đóng vai trò như một công cụ hỗ trợ phân tích ban đầu. Bằng cách tổng hợp nhiều tín hiệu khác nhau, SourceVerify giúp người dùng hiểu vì sao một nội dung có thể đáng nghi, tín hiệu nào ủng hộ giả thuyết AI và tín hiệu nào ủng hộ giả thuyết nội dung thật.
